\chapter{Introduction} \label{chap:intro}

\section{Context and Motivations} \label{sec:context}

For the last couple of years, the field of \ac{ML}, and, in 
particular, its \ac{DL} branch has been a constant presence in the 
computational research breakthroughs, each time creating new mechanisms to 
detect more patterns and allowing the extraction of knowledge from seemingly 
insignificant data, introducing new ways of understanding the world around us.
\cite{tml-enabled-frugal-smart-objects} 

In a similar fashion, the \ac{IoT} paradigm has been gaining more and more
traction, making use of a large number of interconnected devices that gather 
extensive data from the environment and provide a wide range of 
services to users. We can find \ac{IoT} devices in our homes, workplaces, 
and cities, where they continuously monitor and process information to perform 
their designated tasks.

\ac{DNN} models can be used in a wide range of applications, from
automotive and healthcare to smart cities and industry 4.0. \cite{doyu::tinymlaas}
By combining these two paradigms, we can extract even more benefits, as
\ac{ML} models deployed on \ac{IoT} devices enable faster, more
secure and more accurate data processing and analysis. \cite{david::tflm}

\section{Problem} \label{sec:problem}

To provide these tools to users, the most common approach is to
use cloud computing infrastructures, where the data is sent to a
remote server for processing and the results are sent back to the user. However,
this approach presents some limitations, such as latency, privacy and security
concerns, and the need for a constant internet connection.

To overcome these limitations, the concept of edge computing has been
introduced, where the data is processed closer to the source, on edge devices,
such as \acp{MCU}. However, these devices have limited computational
resources, which can make it challenging to deploy complex \ac{DNN} models on
them. \cite{teerapittayannon::ddnn}

There have been some attempts to address this problem, such as NAMP \cite{rego::namp},
which proposes a method for partitioning \ac{DNN} models between \acp{MCU}, relying
on frameworks such as \ac{TFLM} \cite{david::tflm}. However, this approach only allows
for the deployment of relatively lightweight models, such as \textit{LeNet-5} \cite{lenet-5},
and still requires improvement in its scalability potential and communication
protocols.

\section{Hypothesis} \label{sec:hypothesis}

Considering the context and the problem presented, this work aims to answer the
following hypothesis:

\begin{quote}
  ``Is it possible to create a reliable tool to support the deployment of complex DNN models on far edge and edge devices?'' 
\end{quote}

\section{Research Questions} \label{sec:rquestions}

As the hypothesis tackles a complex problem, there are several research
questions that can be derived from it as intermediate steps to answer it,
further completing the understanding of the problem and guiding the
research process.

\begin{itemize}
  \item \textbf{RQ1:} Is it feasible to integrate edge devices into far-edge devices' inference networks without compromising performance?
  \item \textbf{RQ2:} Is it feasible to integrate edge devices into far-edge devices' inference networks without compromising security?
  \item \textbf{RQ3:} Is there a feasible method to fragment more complex \ac{DNN} models (e.g., \acp{ViT}) in the context of edge and far-edge computing?
\end{itemize}

% \section{Dissertation Structure} \label{sec:structure}

% In addition to the introduction, this dissertation contains x more 
% chapters.

% In Chapter~\ref{chap:ch2}, the state of the art is described and 
% related work is presented. 
% Chapter~\ref{chap:ch3}, ipsum dolor sit amet, consectetuer
% adipiscing elit.
